\section{Введение}

Настоящий документ является спецификацией требований к программному обеспечению (СРД) веб-сервиса автоматизированной проверки студенческих отчётов и шаблонов выпускных квалификационных работ (ВКР). Основные пользователи системы~--- студенты: они самостоятельно загружают материалы и получают отчёты с результатами формальной проверки (предупреждения, ошибки).

\subsection{Цель документа}

Цель документа~--- описать функциональные и нефункциональные требования к системе, которая автоматизирует формальную проверку двух типов документов в формате PDF:
\begin{itemize}
  \item шаблон ВКР;
  \item отчёт по учебной практике.
\end{itemize}
Проверка выполняется по формальным критериям (оформление, структура, нумерация, шрифты и т.\,д.); оценка содержания и качества текста не входит в функциональность, которую предполагает система.

\subsection{Проблема}

Студенты часто сдают отчёты и шаблоны с формальными ошибками. Сервис автоматизирует формальную проверку: система формирует и передаёт пользователю отчёт с результатами проверки (предупреждения, ошибки).

\subsection{Область применения}

Веб-сервис предназначен для загрузки PDF-документов студентами, выбора типа документа (шаблон ВКР или отчёт по практике), запуска автоматической формальной проверки и получения отчёта с результатами (предупреждения, ошибки). Поддерживаются загрузка одного или нескольких файлов и просмотр/экспорт отчёта о проверке. Проверяются только документы в формате PDF.

\subsection{Аудитория документа}

Документ адресован разработчикам и заказчикам проекта.
